% Page 041

\seite{41}

\abschnitt{Revision}
\pstart
Am 12.\ Mai wurde die hiesige Schule durch den\ob
Herrn Kreisschulinspektor Bentz aus Bitburg\ob
revidiert.

\pend
\abschnitt{Kirchenfestlichkeiten}
\pstart

Am 3.\ oder hat die Schulen von Malbergweich und\ob
von Sefferweich in der Kirche zu Sefferweich die\ob
gemeinschaftliche Feier des Fronleichnamsfestes durch\ob
einen Gesang mit Gebet\unsicher{} gefeiert, wonach\ob
Verschiedene schöne Vorträge und Lieder verschönerten\ob
den Nachmittag des Festes; wacht in dem jung\ob
Deutscher Geister erster dann auf feine\ob
Anmeldesachen Aufführung die Kaiserwecken durch den\ob
Ortsschulinspector die Kreutzprozession durch besondere\ob
Predigt zum Verständniß warum denn Doch auf\ob
gröbste gegenüber Gemeinschaft.

\pend
\abschnitt{Vertretung}
\pstart

Vom 1.\ oder bis 1.\ August übernahm der hiesige\ob
Lehrer die verwaiste Schuldstelle in Malbergweich.\ob
Der Unterricht wurde abwechselnd einmal morgens\ob
und einmal nachmittags abgehalten.

\pend
\abschnitt{Einquartierung}
\pstart

Am 15.\ September wurde die hiesige Ort mit\ob
einer Theil mit dem Orte eine Eskadron\ob
Dragoner belegt. Am 18.\ August erhielt der\ob
Ort die 11.\ Verpflegungsabteilung auf zwei\ob
Tage in 1 Nacht.

\pend
\abschnitt{Ernte}
\pstart

An Halmfrüchten war die diesjährige Ernte\ob
ziemlich schlecht. Heu, Grummet, Abfall bei\ob
großem Futterbedürfniß. Das Getreide\ob
gut, ob etwas in geringerer Menge. An der\ob
Quantität wurde auch hinlänglich geerntet. Wenn\ob
auch Hegel stehen die Hopfen von Seffern und\ob
Schleidt. Beim andauernden Regen ist zu\ob
viel an Kartoffelfäule liefern die Fang.
\pend
